% =============================================================================
% APPENDIX A: ARTIFACT DESCRIPTION AND EVALUATION INSTRUCTIONS
% =============================================================================
\section{Artifact Description and Evaluation Instructions}
\label{app:artifact}

The submitted folder contains this report, \texttt{declaration.txt}, a \texttt{README.md} with the same content as this appendix in more detail, the LaTeX source of the report, and the \texttt{artifact/} directory described below.

\subsection{Directory Structure}
\begin{table}[H]
\centering
\small
\begin{tabularx}{\textwidth}{l X}
\toprule
\textbf{Path} & \textbf{Contents} \\
\midrule
\texttt{artifact/gem5\_patches/} & \texttt{tlb.\{cc,hh\}.<d>.fixed} for $d \in$ \{inc, excl, nine, agile, perf, colt\}, \texttt{tlb.\{cc,hh\}.perf.strict}, \texttt{pagetable\_walker.cc.colt}, the generators \texttt{make\_colt\_fixed.py} and \texttt{make\_perf\_strict.py}, the checker scripts \texttt{apply\_checker*.py}, diffs against the mid-review snapshots, and \texttt{midreview/} with the mid-review snapshots. \\
\texttt{artifact/benchmarks/} & \texttt{tests/} (tlbstress, pf\_patterns, perf\_holes, gups), \texttt{microbench/} (mid-review microbenchmarks), \texttt{stdbench/} (memcached and Redis drivers), \texttt{spec/} (SPEC CPU2017 build configuration). \\
\texttt{artifact/configs/} & \texttt{fs\_run.py}, the full-system configuration. \\
\texttt{artifact/scripts/} & Build, run, collect and plot automation (Table~\ref{tab:art_scripts}). \\
\texttt{artifact/results/} & \texttt{windows.csv} (554 timing windows), \texttt{spec\_fs\_matrix.csv} (238 SPEC verdicts), \texttt{se\_bigmem.csv}, \texttt{colt\_se\_verification.csv}, \texttt{figures/} with \texttt{numbers.txt}, and \texttt{design\_reports/} (audit and SE reports of each design). \\
\bottomrule
\end{tabularx}
\end{table}

\begin{table}[H]
\centering
\small
\caption{Main scripts.}
\label{tab:art_scripts}
\begin{tabularx}{\textwidth}{l X}
\toprule
\textbf{Script} & \textbf{Role} \\
\midrule
\texttt{build\_design.sh <gem5> <d> [--checker]} & Copies the design's files into a gem5 25.1.0.1 tree (and the walker for CoLT-FA), optionally adds the checker, and builds \texttt{gem5.opt}. \\
\texttt{quick\_start.sh <gem5>} & Five-minute SE check of CoLT-FA with the oracle, with the buffer on and off. Also runs the user's own program. \\
\texttt{build\_microbench.sh}, \texttt{build\_stdbench.sh} & Build the test programs and download and build GAPBS, XSBench, STREAM, memcached, mc-crusher and Redis as static binaries. \\
\texttt{run\_spec\_fs\_shards.sh} & Full-system SPEC CPU2017 runs, one \texttt{gem5} per benchmark, with measurement windows. \\
\texttt{run\_stdbench\_fs3.sh}, \texttt{run\_micro\_fs.sh} & Full-system runs of the standard workloads and microbenchmarks. \\
\texttt{se\_bigmem.sh} & SE runs with 8 GiB and byte-by-byte output comparison. \\
\texttt{gen\_se\_jobs.py}, \texttt{runner.py} & SE correctness job matrix and parallel runner. \\
\texttt{provision\_fs\_host.sh}, \texttt{sync\_fs\_results.sh} & Copy a design's inputs to a new machine; copy results back. \\
\texttt{collect\_results.py}, \texttt{plot\_results.py} & Turn all output into CSV tables and draw every figure of Part~\ref{part:post}. \\
\bottomrule
\end{tabularx}
\end{table}

\subsection{Features and Test Scenarios}
\begin{table}[H]
\centering
\small
\begin{tabularx}{\textwidth}{l X X X}
\toprule
\textbf{Feature} & \textbf{Test scenario} & \textbf{Script} & \textbf{Expected outcome} \\
\midrule
L2 invalidation on unmap and CR3 switch & FS SPEC CPU2017; SE \texttt{tlbstress remap} & \texttt{run\_spec\_fs\_shards.sh}, \texttt{runner.py} & SPEC output check passes; 0 oracle mismatches \\
PCID-aware L2 & SE multi-process jobs; FS runs & \texttt{gen\_se\_jobs.py} & 0 oracle mismatches \\
Runtime parameters & \texttt{NINE\_L2\_SETS}, \texttt{EXCL\_SPILL}, \texttt{AGILE\_PF\_*}, \texttt{COLT\_FA\_ENTRIES}, \texttt{PERF\_STRICT} & any run & statistics change, correctness unchanged \\
Translation oracle and checkers & \texttt{TLBCHK\_FS\_PERIOD}, \texttt{TLBCHK\_MAXWARN}, \texttt{TLBCHK\_PANIC} & \texttt{build\_design.sh --checker} & \texttt{chkOracleMismatch} = 0 \\
CoLT-FA in full system & \texttt{COLT\_ENABLE=1/0}, 8 or 32 entries, huge pages on/off & \texttt{run\_stdbench\_fs3.sh}, \texttt{quick\_start.sh} & CoLT hits only when enabled; fewer walks on graph workloads \\
Measurement windows & \texttt{WINDOW\_WARM}, \texttt{WINDOW\_INSTS}, \texttt{WINDOW\_EXIT} & \texttt{run\_spec\_fs\_shards.sh} & one statistics dump per benchmark \\
\bottomrule
\end{tabularx}
\end{table}

\noindent \textbf{Findings while running the tests.} perlbench fails on every design in full-system mode (gem5 x87 emulation bug); 631.deepsjeng\_s fails in full-system mode (memory limit) and passes in SE mode; 619.lbm\_s triggers the \texttt{gem5} assertion \texttt{Packet::setCacheResponding(): !flags.isSet(CACHE\_RESPONDING)} on every design right after its window, and passes in SE mode; the first Redis runs made no progress inside the guest. The SE-mode failures (\texttt{setitimer}, fork with shared memory, multi-process remap) are identical on stock \texttt{gem5}.

\subsection{Assumptions and Unsupported Features}
\texttt{gem5} 25.1.0.1 for x86, gcc-11, Python 3.10 and scons 4.8.1. Full-system runs need the \texttt{gem5} resources kernel \texttt{vmlinux-x86-ubuntu-6.8.0-52-generic} and disk image \texttt{x86-ubuntu-24.04.img}, and a licensed SPEC CPU2017 installation (not included). Single simulated core and at most 3 GiB guest memory. Without KVM the guest boots on the atomic CPU. The Hybrid L2 TLB is included only as a mid-review snapshot.

\subsection{Getting Started (about 30 minutes with a built tree)}
\begin{table}[H]
\centering
\small
\begin{tabularx}{\textwidth}{l X l l}
\toprule
\textbf{Step} & \textbf{Command} & \textbf{Human} & \textbf{Compute} \\
\midrule
1 & \texttt{git clone https://github.com/gem5/gem5 -b v25.1.0.1} & 2 min & 5 min \\
2 & \texttt{bash artifact/scripts/build\_design.sh <gem5> colt --checker -j 16} & 1 min & 1--2 h clean, 6 min on a built tree \\
3 & \texttt{bash artifact/scripts/quick\_start.sh <gem5>} & 1 min & under 1 min \\
4 & \texttt{python3 artifact/scripts/plot\_results.py artifact/results /tmp/fig} & 1 min & 10 s \\
\bottomrule
\end{tabularx}
\end{table}

\noindent In our test on a stock \texttt{gem5} 25.1.0.1 tree, step 2 took 6 minutes and step 3 took under a minute. For the ascending pattern it reported 14,330 CoLT hits and 10,329 walks with CoLT on, against 24,659 walks with CoLT off. Step 3 should print \texttt{oracleMismatch=0} on every line, and CoLT hits above zero with fewer walks when CoLT is enabled. To use one's own program: \texttt{QS\_PROG=/path/to/static\_binary QS\_ARGS="..." bash artifact/scripts/quick\_start.sh <gem5>}.

\subsection{Detailed Evaluation}
\begin{table}[H]
\centering
\small
\begin{tabularx}{\textwidth}{l X X l X}
\toprule
\textbf{\#} & \textbf{Experiment and purpose} & \textbf{How to run} & \textbf{Runtime} & \textbf{Expected result and location} \\
\midrule
E1 & SE correctness matrix of a design against stock \texttt{gem5} & \texttt{gen\_se\_jobs.py}, then \texttt{runner.py <jobs> --jobs 24} & 3--6 h & 0 mismatches; Table~\ref{tab:se_matrix}; \texttt{design\_reports/} \\
E2 & CoLT-FA SE checks, buffer on/off & \texttt{quick\_start.sh} or the 36-run loop & 10 min & 36/36 PASS; Table~\ref{tab:colt_se}; \texttt{colt\_se\_verification.csv} \\
E3 & SPEC CPU2017 full system, test inputs & \texttt{run\_spec\_fs\_shards.sh <d> 43 atomic} with \texttt{WINDOW\_INSTS=200000000 WINDOW\_WARM=50000000} & 1--3 days & 37/43 PASS; Table~\ref{tab:spec_outcome}; \texttt{spec\_fs\_matrix.csv} \\
E4 & SE 8 GiB runs of deepsjeng\_s and lbm\_s & \texttt{se\_bigmem.sh <bench> inc excl nine agile perf colt} & 6--16 h & PASS on all; Table~\ref{tab:se_bigmem}; \texttt{se\_bigmem.csv} \\
E5 & TLB-intensive workloads per design, huge pages on/off & \texttt{build\_stdbench.sh}; \texttt{[THP=never] [COLT\_ENABLE=0] STD\_TAG=<t> run\_stdbench\_fs3.sh <d>} & 4--7 h & Tables~\ref{tab:thpoff_all} and~\ref{tab:colt_vs_nine}; \texttt{windows.csv} \\
E6 & SPEC reference windows & \texttt{SPEC\_FS\_SIZE=ref WINDOW\_WARM=3000000000 WINDOW\_EXIT=1 run\_spec\_fs\_shards.sh <d> 9 atomic <benches>} & about 5 h & Table~\ref{tab:ref}; \texttt{windows.csv} \\
E7 & Perforated Pages strict mode & build \texttt{perfstrict}, run E5/E6 with \texttt{PERF\_STRICT=1} & as E5/E6 & Table~\ref{tab:perf_strict} \\
E8 & Tables and plots & \texttt{collect\_results.py <mirror> out}; \texttt{plot\_results.py out out/figures} & 1 min & all figures of Part~\ref{part:post} \\
\bottomrule
\end{tabularx}
\end{table}
